\abstract

This thesis investigates whether tool-orchestrated agentic workflows can improve reverse engineering and malware analysis outcomes relative to weaker baselines. The system under study does not rely on a single monolithic agent. Instead, it coordinates planners, specialist workers, validators, and reporters through MCP-style tool access to Ghidra, static-analysis utilities, dynamic-analysis infrastructure, and supporting evidence-recovery tools. The thesis focuses on how staged reasoning, explicit tool routing, and evidence review affect tasks such as behavior reconstruction, configuration decoding, packing analysis, and malware-oriented triage.

This draft abstract is intentionally conservative. It establishes the thesis scope and argument structure without claiming finished results. \TODO{Revise the abstract after the architecture, evaluation, and final quantitative findings are stable.} \TODO{State the primary research questions and final contribution claims in two or three sentences.}
